\clearpage\pagestyle{otMain}\setlength{\otBlocEspacement}{2.5mm}\fontsize{11}{15}\selectfont\corrigesChapitre{5}{Les fractions comme nombres}
\begin{blocLivre}{cor-5-1}\begin{corrige}{1}{Quotient exact}\demarche{Le quotient exact s’écrit avec une fraction.}
$\res{\frac58}$ ; $\res{\frac{13}4}$ ; $\res{\frac27}$. $\frac58=0{,}625$ et $\frac{13}4=3{,}25$ sont décimaux ; la division de 2 par 7 ne s’arrête pas : $\frac27$ n’est pas décimal.
\end{corrige}\end{blocLivre}
\begin{blocLivre}{cor-5-2}\begin{corrige}{2}{Égalités à trous}\demarche{Le nombre qui multiplié par b donne a est a sur b.}
$\square=\res{\frac49}$ ; $\square=\res{\frac35}$ ; $\square=\res1$ ; $\square=\res{\frac{18}6}=3$.
\end{corrige}\end{blocLivre}
\begin{blocLivre}{cor-5-3}\begin{corrige}{3}{À connaître}\demarche{Automatismes du programme.}
$\res{0{,}25}$ ; $\res{0{,}75}$ ; $\res{1{,}5}$ ; $\res{2{,}5}$ ; $\res2$. Et $0{,}5=\res{\frac12}$.
\end{corrige}\end{blocLivre}
\begin{blocLivre}{cor-5-4}\begin{corrige}{4}{Lire sur la droite}\demarche{Compter les pas d’un quart depuis 0.}
A : $\res{\frac34}$. B : $\res{\frac74}=1+\frac34$. C : $\res{\frac{10}4}=2+\frac24$, soit $2+\frac12$.
\label{fig-cor-5-4}\figureGeo[corrige]{fq-4}
\end{corrige}\end{blocLivre}
\begin{blocLivre}{cor-5-5}\begin{corrige}{5}{Le quart de 3}\demarche{Une fraction est le quotient du numérateur par le dénominateur.}
Une part mesure $3:4=\res{\frac34}$ d’unité : quatre parts de $\frac34$ font bien $4\times\frac34=3$. C’est exactement trois quarts d’une unité.
\label{fig-cor-5-5}\figureGeo[corrige]{fq-5}
\end{corrige}\end{blocLivre}
\begin{blocLivre}{cor-5-6}\begin{corrige}{6}{Placer des fractions}\demarche{Écrire chaque fraction en sixièmes.}
D : 5 pas d’un sixième. E : 11 pas, soit $1+\frac56$. F : $\frac43=\frac86$, soit 8 pas, $1+\frac26$. Les points figurent en vert.
\label{fig-cor-5-6}\figureGeo[corrige]{fq-6}
\end{corrige}\end{blocLivre}
\begin{blocLivre}{cor-5-7}\begin{corrige}{7}{Entier, décimal ou non décimal}\demarche{Effectuer la division du numérateur par le dénominateur.}
Entier : $\res{\frac{15}5=3}$. Décimaux non entiers : $\res{\frac74=1{,}75}$ et $\res{\frac9{10}=0{,}9}$. Non décimaux : $\res{\frac23}$ et $\res{\frac56}$, car les divisions ne s’arrêtent pas.
\end{corrige}\end{blocLivre}
\begin{blocLivre}{cor-5-8}\begin{corrige}{8}{Nombres mixtes}\demarche{Extraire les unités entières.}
$\frac{17}5=\res{3+\frac25}$ ; $\frac{23}6=\res{3+\frac56}$ ; $3+\frac27=\frac{21}7+\frac27=\res{\frac{23}7}$.
\end{corrige}\end{blocLivre}
\begin{blocLivre}{cor-5-9}\begin{corrige}{9}{Graduer un segment}\demarche{Utiliser le guide-âne plutôt que la règle graduée.}
Les droites du réseau coupent le segment [AB] en trois parts égales, marquées en vert. Une part mesure $\res{\frac53}$ cm, soit environ $\res{1{,}67}$ cm.
\label{fig-cor-5-9}\figureGeo[corrige]{fq-9}
\end{corrige}\end{blocLivre}
\begin{blocLivre}{cor-5-10}\begin{corrige}{10}{Encadrer entre deux entiers}\demarche{Encadrer le numérateur par des multiples du dénominateur.}
$15<17<18$ : $\res{5<\frac{17}3<6}$. $\res{4<\frac{47}{10}<5}$. $21<22<28$ : $\res{3<\frac{22}7<4}$.
\end{corrige}\end{blocLivre}
\begin{blocLivre}{cor-5-11}\begin{corrige}{11}{Le partage des pizzas}\demarche{La part est le quotient du nombre de pizzas par le nombre d’enfants.}
Première table : $3:4=\res{\frac34}$ de pizza. Seconde : $5:8=\res{\frac58}$. Or $\frac34=\frac68>\frac58$ : on mange plus à la $\res{\text{première table}}$.
\end{corrige}\end{blocLivre}
\begin{blocLivre}{cor-5-12}\begin{corrige}{12}{La ficelle}\demarche{Distinguer valeur exacte et valeur approchée.}
$\frac25=\res{0{,}4}$ m, valeur exacte. $\frac23$ m n’est pas décimal : environ $\res{0{,}67}$ m.
\end{corrige}\end{blocLivre}
\begin{blocLivre}{cor-5-13}\begin{corrige}{13}{Le nombre caché}\demarche{Utiliser la définition du quotient.}
C’est $\frac68$, qui se simplifie en $\res{\frac34}$, soit $\res{0{,}75}$. Contrôle : $8\times0{,}75=6$.
\end{corrige}\end{blocLivre}
\begin{blocLivre}{cor-5-14}\begin{corrige}{14}{Les nombres rompus}\demarche{Le partage exact s’écrit avec une fraction.}
Chacun reçoit $\res{\frac73}$ d’arpent, soit $\res{2+\frac13}$ arpents. Comme $\frac13=0{,}333\ldots$, l’arrondi au dixième est $\res{2{,}3}$.
\end{corrige}\end{blocLivre}
\begin{blocLivre}{cor-5-15}\begin{corrige}{15}{Vrai ou faux}\demarche{Utiliser la définition : 3 fois un tiers égale 1.}
$\res{\text{Non}}$ : $3\times0{,}33=0{,}99$, qui n’est pas égal à 1. $0{,}33$ n’est qu’une valeur approchée de $\frac13$ ; on écrit $\frac13\approx0{,}33$.
\end{corrige}\end{blocLivre}
